% Copyright (c) 2026 Angela Villota, and collaborators from the CyED block
% Licensed under the PolyForm Noncommercial License 1.0.0.
% Commercial use is prohibited without prior written authorization.
%===============================================================================
% LaTeX  Diapositivas Computaci�n y Estructuras Discretas III 2023-2 ICESI
% Andr��s A. Aristiz�bal P. - �ngela P. Villota G.
%===============================================================================
%===============================================================================
\documentclass[compress,xcolor={dvipsnames}]{beamer}
\mode<presentation> {
%============================================================================== 
\usetheme{Boadilla}
%\usetheme{Copenhagen}
%\usetheme{CambridgeUS}
%============================================================================== 
%\usecolortheme{crane}
%\usecolortheme{seahorse}
%============================================================================== 
%\usecolortheme{crane}
\usecolortheme{whale}
%============================================================================== 
\usefonttheme[onlylarge]{structurebold}

%==============================================================================
%==============================================================================
\setbeamertemplate{navigation symbols}{}
\setbeamertemplate{footline}[frame number]
\setbeamerfont*{frametitle}{size=\normalsize,series=\bfseries}
%\setbeamercolor{structure}{fg=black!80!white,bg=white}
%\setbeamercolor{palette quaternary}{fg=white!80!gray,bg=red!50!black}}

\setbeamercolor{structure}{fg=white,fg=blue!60!black}  %%%%Estaba en 70
\setbeamercolor{palette quaternary}{fg=gray!40!black,bg=gray!20!white}}


%=============================================================================== 
%========================================================================================
\usepackage[spanish,activeacute]{babel}
\usepackage[applemac]{inputenc}
\usepackage[spanish]{babel}
\usepackage{times}
\usepackage{latexsym}
\usepackage{color}
\usepackage{amsmath}
\usepackage{amsfonts}
\usepackage{amssymb}
\usepackage{stmaryrd}
\usepackage{indentfirst}
\usepackage{multicol}
\usepackage{graphicx}
\usepackage{graphics}
\usepackage{newlfont}
\usepackage{exscale}
\usepackage[T1]{fontenc}

\usepackage{tikz}
\usepackage{amscd}
\usepackage{pb-diagram}
%\usepackage{pstricks}
\usepackage{makeidx}
\usepackage{textpos}

%==============================================================================
%==============================================================================
\tikzstyle{block}=[draw opacity=0.7,line width=1.4cm]
%============================================================================== 
%=======================================================================================
\newtheorem{teo}{Teorema}
\newtheorem{prop}{Proposition}
\newtheorem{afir}{Afirmaci�n}
\newtheorem{obs}{Observaci�n}
\newtheorem{lem}{Lemma}
\newtheorem{cor}{Corollary}
\newtheorem{ej}{Example}
\newtheorem{ejer}{Exercise}
\newtheorem{ejers}{Exercises}
\newtheorem{ejs}{Examples}
\newtheorem{df}{Definition}

%%%%%%%%%%%%%%%%%%%%%%%%%%%%%%%%%%%%%%%%%%%%%%%%%%%%%%%%%%%%%%%%%%%%%Flechas
\newcommand{\stateequiv}{\,\sim\,}
\newcommand{\entonces}{\,\rightarrow\,}
\newcommand{\pertenece}{\,\in \,}
\newcommand{\talque}{\,| \,}
\newcommand{\sii}{\,\leftrightarrow\,}
\newcommand{\ra}{\,\rightarrow\,}
\newcommand{\lra}{\leftrightarrow}
\newcommand{\Ra}{\;\Longrightarrow\;}
\newcommand{\lora}{\longrightarrow}
\newcommand{\we}{\,\wedge\,}
\newcommand{\ve}{\,\vee \,}
\newcommand{\bbsq}{\blacksquare\!\blacksquare}
\newcommand{\powerset}{\raisebox{.15\baselineskip}{\Large\ensuremath{\wp}}}
%%%%%%%%%%%%%%%%%%%%%%%%%%%%%%%%%%%%%%%%%%%%%%%%%%%%%%%%%%%%%%%%%%%%%S�mbolos l�gicos
\newcommand{\mo}{\models}
\newcommand{\lan}{\langle}
\newcommand{\ran}{\rangle}
\newcommand{\ap}{\approx}
%%%%%%%%%%%%%%%%%%%%%%%%%%%%%%%%%%%%%%%%%%%%%%%%%%%%%%%%%%%%%%%%%%%%%Letras
\newcommand{\vep}{\varepsilon}
\newcommand{\vph}{\varphi}
\newcommand{\de}{\delta}
\newcommand{\si}{\sigma}
\newcommand{\Si}{\Sigma}
\newcommand{\om}{\omega}
\newcommand{\al}{\alpha}
\newcommand{\be}{\beta}
\newcommand{\vom}{\varOmega}
\newcommand{\vTh}{\varTheta}
\newcommand{\vth}{\vartheta}
\newcommand{\vDe}{\varDelta}
\newcommand{\wh}{\widehat}
\newcommand{\nn}{\~n}
\newcommand{\m}{i}
\newcommand{\ps}{P_{\Sigma}}
\newcommand{\psone}{P_{\Sigma^{*}}}
\newcommand{\psoneast}{P_{\Sigma'}}
\newcommand{\ts}{T_{\Sigma}}
\newcommand{\dis}{\displaystyle}
\newenvironment{variableblock}[3]{%
  \setbeamercolor{block body}{#2}
  \setbeamercolor{block title}{#3}
  \begin{block}{#1}}{\end{block}}
  
\usepackage{listings}
\usepackage{color}

\definecolor{dkgreen}{rgb}{0,0.6,0}
\definecolor{gray}{rgb}{0.5,0.5,0.5}
\definecolor{mauve}{rgb}{0.58,0,0.82}
  
  \lstset{frame=tb,
  language=Python,
  aboveskip=3mm,
  belowskip=3mm,
  showstringspaces=false,
  columns=flexible,
  basicstyle={\small\ttfamily},
  numbers=none,
  numberstyle=\tiny\color{gray},
  keywordstyle=\color{blue},
  commentstyle=\color{dkgreen},
  stringstyle=\color{mauve},
  breaklines=true,
  breakatwhitespace=true,
  tabsize=3
}

%============================================================================== 
%==========================================================================================

\title{Computaci\'on y Estructuras Discretas III}

\author[Andr�s A. Aristiz�bal P.]{Andr�s A. Aristiz�bal P. aaaristizabal@icesi.edu.co \\  Angela Villota - apvillota@icesi.edu.co}

\institute[Universidad ICESI]{Departamento de Computaci\'on y Sistemas Inteligentes \\
\includegraphics[height=1.5cm]{../logoICESI2.pdf}}

\date[]{\today}

%==============================================================================
\AtBeginSubsection[]{

%%%%%%%%%%%%%%%%%%%%%%%%%%%%%%%%%%%%%%%%%%%%%%%%%%%%%%%%%%%%%%%%%%%%%%%%%%%%%%
\begin{frame}
    \frametitle{Agenda del d�a}
    \tableofcontents[currentsubsection, 
    hideothersubsections, 
    sectionstyle=show/shaded, 
    subsectionstyle=show/shaded,]
  \end{frame}}
%%%%%%%%%%%%%%%%%%%%%%%%%%%%%%%%%%%%%%%%%%%%%%%%%%%%%%%%%%%%%%%%%%%%%%%%%%%%%%

%============================================================================== CUERPO DEL DOCUMENTO
%%%%%%%%%%%%%%%%%%%%%%%%%%%%%%%%%%%%%%%%%%%%%%%%%%%%%%%%%%%%%%%%%%%%%%%%%%%%%%
\begin{document}
\maketitle
\setlength{\parindent}{0cm}
%%%%%%%%%%%%%%%%%%%%%%%%%%%%%%%%%%%%%%%%%%%%%%%%%%%%%%%%%%%%%%%%%%%%%%%%%%%%%%
%\begin{frame}
 %   \titlepage
%\end{frame}
%%%%%%%%%%%%%%%%%%%%%%%%%%%%%%%%%%%%%%%%%%%%%%%%%%%%%%%%%%%%%%%%%%%%%%%%%%%%%%

%%%%%%%%%%%%%%%%%%%%%%%%%%%%%%%%%%%%%%%%%%%%%%%%%%%%%%%%%%%%%%%%%%%%%%%%%%%%%%
\begin{frame}
\frametitle{Agenda del d�a}
  \tableofcontents
\end{frame}
%%%%%%%%%%%%%%%%%%%%%%%%%%%%%%%%%%%%%%%%%%%%%%%%%%%%%%%%%%%%%%%%%%%%%%%%%%%%%%

%%%%%%%%%%%%%%%%%%%%%%%%%%%%%%%%%%%%%%%%%%%%%%%%%%%%%%%%%%%%%%%%%%%%%%%%%%%%%%

\setbeamertemplate{itemize/enumerate body begin}{\footnotesize}
\setbeamertemplate{itemize/enumerate subbody begin}{\footnotesize}

\section[Regular languages and Automata theory]{Regular languages and Automata theory}

\subsection{Deterministic Automata}
{\footnotesize \begin{frame}{Regular languages and Automata theory}
What is a finite automaton? \pause
\begin{block}{}
A finite automaton is an abstract machines that process input \textbf{strings}, which are either accepted or rejected. \pause
\begin{center}
\includegraphics[width=8cm]{../AFD} 
\end{center}
\end{block}
\end{frame}}

{\footnotesize \begin{frame}{Regular languages and Automata theory}
How does a finite automaton work? \pause
\begin{block}{}
The automaton operates by reading the symbols written on a semi-infinite tape divided into cells, onto which an input string $u$ is written, one symbol per cell. \pause
\end{block}
What are the components of a finite automaton? \pause
\begin{block}{}
The automaton has a control unit or reading head, which has a finite number of internal configurations, called states. \pause There are some states of particular importance: these are the initial state and the final or accepting states.
\end{block}
\end{frame}}

{\footnotesize \begin{frame}{Regular languages and Automata theory}
What is the formal definition of a finite automaton? \pause
\begin{df}
A finite automaton $M$ is defined as a 5-tuple, such that $M =(\Sigma,Q,q_0,F,\delta)$. \pause
\begin{enumerate}
\item An alphabet $\Sigma$, called the tape alphabet. All strings processed by $M$ belong to $\Sigma^*$. \pause
\item $Q={q_0,q_1,...,q_n}$, the set of internal states of the automaton. \pause
\item$q_0 \in Q$, the initial state. \pause
\item $F \subseteq Q$, the set of final or accepting states. $F \neq \emptyset$. \pause
\item The transition function of the automaton:
\begin{center}
\begin{tabular}{rcc}
$\delta \, : \, Q \times \Sigma$ & $\longrightarrow$ & Q \\
$(q,s)$ & $\longmapsto$ & $\delta(q,s)$
\end{tabular}
\end{center}
\end{enumerate}
\end{df}
\end{frame}}

{\footnotesize \begin{frame}{Regular languages and Automata theory}
How does a finite automaton start working? \pause
\begin{block}{}
An input string $u$ is placed on the tape in such a way that the first symbol of $u$ occupies the first cell of the tape. The reading head is initially in the initial state $q_0$. \pause
\begin{center}
\includegraphics[width=9cm]{../AFD1a} 
\end{center}
\end{block}
\end{frame}}

{\scriptsize \begin{frame}{Regular languages and Automata theory}
How does a finite automaton transition from state to state? \pause
\begin{block}{}
The transition function $\delta$ indicates the state to which the finite control moves, depending on the scanned symbol and its current state. \pause So $\delta(q,s)=q'$, which indicates that when symbol $s$ is encountered, the reading head moves from state $q$ to $q'$, shifting to the right on the tape. \pause This is one computational step. \pause
\begin{center}
\includegraphics[width=9cm]{../AFD2a} \pause 
\end{center}
Since the transition function $\delta$ is defined for every state-symbol combination, an input string is processed entirely until the control unit encounters the empty string.
\end{block}
\end{frame}}

{\footnotesize \begin{frame}{Regular languages and Automata theory}
\begin{ej}
Let's consider the automaton defined by the following $5$ components:
\begin{tabular}{l}
$\Sigma = {a,b}$ \\
$Q = {q_0,q_1,q_2}$ \\
$q_0$ : initial state. \\
$F={q_0,q_2}$, accepting states. \\
Transition function $\delta:$ \pause
\end{tabular}
\begin{center}
\includegraphics[width=9cm]{../AFD3a}
\end{center}
\end{ej}
\end{frame}}

{\scriptsize \begin{frame}{Regular languages and Automata theory}
\begin{ej}
We illustrate the processing of two input strings. \pause
\begin{enumerate}
\item $u=aabab$ \pause
\begin{center}
\includegraphics[width=6cm]{../AFD4a} \pause
\end{center}
Since $q_2$ is an accepting state, the input string $u$ is accepted. \pause
\item $v=aababa$ \pause
\begin{center}
\includegraphics[width=6cm]{../AFD5a} \pause
\end{center}
Since $q_1$ is not an accepting state, the string $v$ is rejected.
\end{enumerate}
\end{ej}
\end{frame}}

{\footnotesize \begin{frame}{Regular languages and Automata theory}
How do we call these previously defined automata? \pause
\begin{block}{}
They are called Deterministic Finite Automata (DFA) because for each state $q$ and for each symbol $a \in \Sigma$, the transition function $\delta(q,a)$ is always defined. \pause
\end{block}
How do we call the language accepted by a finite automaton? \pause
\begin{block}{}
Let $M$ be an automaton, the language accepted by this automaton is denoted as $L(M)$, and it is defined as follows: \pause
\begin{center}
$L(M) := \{u \in \Sigma^* \mid M$ terminates processing the input string $u$ in a state $q \in F \}$.
\end{center}
\end{block}
\end{frame}}

{\footnotesize \begin{frame}{Regular languages and Automata theory}
What is a transition diagram of an automaton? \pause
\begin{block}{}
It is a directed and labeled graph that represents the finite automaton. \pause
\end{block}
How is the transition diagram of an automaton obtained? \pause
\begin{block}{}
\begin{tabular}{l}
\includegraphics[width=6cm]{../paso1} \\ \pause
\includegraphics[width=6cm]{../paso2} \\ \pause
\includegraphics[width=6cm]{../paso3} \\ \pause
\includegraphics[width=6cm]{../paso4} \\ \pause
\includegraphics[width=6cm]{../paso5}
\end{tabular}
\end{block}
\end{frame}}

{\footnotesize \begin{frame}{Regular languages and Automata theory}
\begin{ej}
\begin{center}
\includegraphics[width=7cm]{../AFD6}
\end{center}
\end{ej}
\end{frame}}
%%%%%%%%%%%%%

\subsection{Exercises}

\subsection{Non-deterministic Automatons}
{\scriptsize \begin{frame}{Regular languages and Automata theory}
What is a non-deterministic finite automaton? \pause
\begin{block}{}
The non-deterministic finite automata resemble deterministic finite automata, except that for each state $q \in Q$ and each $a \in \Sigma,$ the transition $\delta(q,a)$ can consist of more than one state or may not be defined. \pause
\end{block}
How would we formally define a non-determinist finite automaton? \pause
\begin{block}{}
An NFA is defined by $M=(\Sigma,Q,q_0,F,\Delta)$ where: \pause
\begin{enumerate}
\item $\Sigma$ is the tape alphabet. 
\item $Q$ is the (finite) set of internal states.
\item $q_0 \pertenece Q$ is the initial state. 
\item $\emptyset \neq F \subseteq Q$ is the set of final or accepting states.
\item Let $\powerset(Q)$ be the set of subsets of $Q,$ 
\begin{center}
\begin{tabular}{rcl}
$\Delta \, : \, Q \times \Sigma$ & $\longrightarrow$ & $\powerset(Q)$ \\
$(q,s)$ & $\mapsto$ & $\Delta(q,s) = \{q_{i1},q_{i2},...,q_{ik}\}$
\end{tabular}
\end{center}
\end{enumerate}
\end{block}
\end{frame}}

{\footnotesize \begin{frame}{Regular languages and Automata theory}
What is the meaning $\Delta(q,s) = \{q_{i1},q_{i2},...,q_{ik}\}$? \pause
\begin{block}{}
Being in state $q$, upon encountering symbol $s$, it can randomly transition to any of the states $q_{i1}, q_{i2}, ..., q_{ik}$ and moves to the right. \pause
\end{block}
What if $\Delta(q,s)=\emptyset$? \pause
\begin{block}{}
This means that if, during the processing of an input string $u$, $M$ enters state $q$ while reading symbol $s$ on the tape, the computation is aborted. 
\begin{center}
\includegraphics[width=5cm]{../AFN1} 
\end{center}
\end{block}
\end{frame}}

{\footnotesize\begin{frame}{Regular languages and Automata theory}
How do we represent an NFA by means of a transition diagram? \pause
\begin{block}{}
It is represented similarly to a DFA, the difference lies in the fact that from the same node (state), it may happen that two or more arcs with the same label emerge. 
\begin{center}
\includegraphics[width=4cm]{../AFN2} 
\end{center}
\end{block}
\end{frame}}


{\footnotesize\begin{frame}{Regular languages and Automata theory}
How does an FNA process an input string? \pause
\begin{block}{}
An NFA can process an input string $u \in \Sigma^*$ in different ways. Seen in the transition diagram, this means that there can be various paths from the initial state, labeled with the symbols of $u$.
\end{block}
\end{frame}}

{\footnotesize\begin{frame}{Regular languages and Automata theory}
What is the notion of acceptance for nondeterministic finite automata? \pause
\begin{block}{}
\begin{center}
\begin{tabular}{lll}
$L(M)$ & $=$ & language accepted or recognized by $M$ \\
& $=$ & $\{u \in \Sigma \mid$ there exists at least one complete computation \\
& & of $u$ that ends in a state $q \in F\}$
\end{tabular}
\end{center} 
For a string $u$ to be accepted, there must exist some computation in which $u$ is processed completely and ends with $M$ in an accepting state.
\end{block}
\end{frame}}

{\footnotesize\begin{frame}{Regular languages and Automata theory}
\begin{ej}
Let $M$ be the following NFA: 
\begin{center}
\includegraphics[width=7cm]{../AFN3} 
\end{center}
For the string $u = abb$, there are computations that lead to rejection, aborted computations, and computations that end in accepting states.
\end{ej}
\end{frame}}

{\tiny\begin{frame}{Regular languages and Automata theory}
\begin{ej}
According to the definition of accepted language, $u \in L(M)$, we have: \pause
\begin{center}
\begin{tabular}{l}
\includegraphics[width=5.3cm]{../AFN4} \\ \pause
\includegraphics[width=5.3cm]{../AFN5} \\ \pause
\includegraphics[width=5.3cm]{../AFN6} \\ \pause
\includegraphics[width=5.3cm]{../AFN7}
\end{tabular}
\end{center}
\end{ej}
\end{frame}}
%%%%%%%%%%%%
{\footnotesize\begin{frame}{Regular languages and Automata theory}
Are DFAs and NFAs computationally equivalent? \pause
\begin{block}{}
The DFAs and NFAs are computationally equivalent. \pause It can be seen that a DFA $M=(\Sigma,Q,q_0,F,\delta)$ can be considered as an NFA $M'= (\Sigma,Q,q_0,F,\Delta)$ by defining $\Delta(q,a)={\delta(q,a)}$ for each $q \in Q$ and each $a \in \Sigma$ \pause
\end{block}
For the converse of this implication, we have the following theorem: \pause
\begin{teo}
Given an NFA $M=(\Sigma,Q,q_0,F,\Delta)$, we can construct an equivalent DFA $M'$ to $M$, meaning that $L(M) = L(M')$
\end{teo}
\end{frame}}

{\footnotesize \begin{frame}{Regular languages and Automata theory}
What is a nondeterministic finite automaton with $\lambda$ transitions? \pause
\begin{block}{}
It is a nondeterministic automaton $M=(\Sigma,Q,q_0,F,\Delta)$ where the transition function is defined as $\Delta \, : \, Q \times (\Sigma \cup \{\lambda\}) \rightarrow \powerset(Q).$ \pause
\end{block}
What is a $\lambda$ transiton? \pause
\begin{block}{}
The $\lambda$ transition, $\Delta(q,\lambda)=\{p_{i_1},...,p_{i_n}\}$, also called null or spontaneous transition, allows transitioning between states without processing any symbol read on the tape. \pause
\end{block}
How do we represent lambda transitions in a transition diagram? \pause
\begin{block}{}
In a transition diagram, these transitions result in arcs labeled with $\lambda$.
\end{block}
\end{frame}}

{\footnotesize\begin{frame}{Regular languages and Automata theory}
When is an input $w$ accepted by an NFA-$\lambda$? \pause
\begin{block}{}
A string $w$ is accepted by an NFA-$\lambda$ if there exists at least one path from the initial state $q_0$ to an accepting state, whose labels consist exactly of the symbols of $w$, interspersed with zero or more $\lambda$s.\pause
\end{block}
\begin{ej}
Let $M$ be the following NFA-$\lambda$: \pause
\begin{center}
\includegraphics[width=3cm]{../AFNL1a} 
\end{center}
\end{ej}
\end{frame}}

{\footnotesize \begin{frame}{Regular languages and Automata theory}
\begin{ej}
The acceptance computation for $u=aab$ is: \pause
\begin{center}
\includegraphics[width=4cm]{../AFNL2a} \pause
\end{center}
The acceptance computation for $v=abaa$ is: \pause
\begin{center}
\includegraphics[width=4cm]{../AFNL3a}
\end{center}
\end{ej}
\end{frame}}

{\scriptsize\begin{frame}{Regular languages and Automata theory}
What advantages does an NFA-$\lambda$ have? \pause
\begin{block}{}
It allows even more freedom in the design of automata, especially when there are numerous concatenations. \pause
\end{block}
\begin{ej}
Let $M$ be the following DFA that accepts the language $a^*b^*c^*$ over $\Sigma=\{a,b,c\}$: \pause
\begin{center}
\includegraphics[width=3cm]{../AFD1} \pause
\end{center}
$M'$ is an NFA-$\lambda$ that accepts $L:$ \pause
\begin{center}
\includegraphics[width=3cm]{../AFNL4a} 
\end{center}
\end{ej}
\end{frame}}

{\footnotesize\begin{frame}{Regular languages and Automata theory}
\begin{ej}
Let $M_1$ be the following DFA that accepts the language of all strings with an even number of 'a's over $\Sigma=\{a,b\}$: \pause
\begin{center}
\includegraphics[width=2cm]{../AFD2} \pause
\end{center}
And let $M_2$ be the following DFA that accepts the language of all strings with an even number of 'b's over $\Sigma=\{a,b\}$: \pause
\begin{center}
\includegraphics[width=2cm]{../AFD3} 
\end{center}
\end{ej}
\end{frame}}

{\footnotesize\begin{frame}{Regular languages and Automata theory}
\begin{ej}
$M_3$ is the following NFA-$\lambda$ that accepts the language of all strings with an even number of 'a's or an even number of 'b's over $\Sigma=\{a,b\}:$ \pause
\begin{center}
\includegraphics[width=4cm]{../AFNL5a} 
\end{center}
\end{ej}
\end{frame}}

{\footnotesize\begin{frame}{Regular languages and Automata theory}
Can infinite loops occur in an NFA-$\lambda$? \pause
\begin{block}{}
Unlike DFAs and NFAs, in NFAs-$\lambda$, there can be computations that never terminate. This can happen if the automaton enters a state from which there are multiple chained $\lambda$ transitions that return to the same state. \pause
\begin{center}
\includegraphics[width=5cm]{../AFNL6a} 
\end{center}
\end{block}
\end{frame}}
%%%%%%%%%%%%
{\footnotesize\begin{frame}{Regular languages and Automata theory}
Are NFAs and NFAs-$\lambda$ computationally equivalent? \pause
\begin{block}{}
NFA-$\lambda$ are computationally equivalent to NFAs. In other words, $\lambda$ transitions can be eliminated by adding transitions that simulate them, without altering the accepted language. \pause \\
An NFA $M=(\Sigma,Q,q_0,F,\Delta)$ can be considered as an NFA-$\lambda$ where there are simply zero $\lambda$ transitions. \pause
\end{block}
For the converse of this implication, we have the following theorem: \pause
\begin{teo}
Given an NFA-$\lambda$ $M=(\Sigma,Q,q_0,F,\Delta)$, we can construct an equivalent NFA $M'$ to $M$, meaning that $L(M) = L(M')$.
\end{teo}
\end{frame}}


\subsection{Exercises}

\end{document}


